\clearpage\begingroup\NoSourceHead\SourceLeaf{519b}
\renewcommand{\thefootnote}{\arabic{footnote})}\setcounter{footnote}{13}
\ArithmeticHead{62}{L}{BOLZA, GAUSS UND DIE VARIATIONSRECHNUNG.}
\noindent\hspace*{\parindent}b) Die Untersuchungen von \textsc{Beltrami} und \textsc{Hilbert}.
\par\noindent\hspace*{\parindent}11. Die Formeln (30) wurden explizite zuerst von \textsc{Beltrami}\footnote{E. \textsc{Beltrami}, \textit{Sulla teoria delle linee geodetiche}, Rendiconti del Reale Istituto Lombardo (II),
Bd. I (1868), S. 708; auch Opere matematiche, Bd. I, S. 366.} in einer
Arbeit aus dem Jahr 1868 angegeben und zusammen mit der partiellen Differentialgleichung
(29) als spezielle Fälle allgemeiner Variationssätze erkannt
und aus letzteren abgeleitet. \textsc{Beltrami} beweist\footnote{\textsc{Beltrami} und ebenso später \textsc{Hilbert} nehmen das Integral $J$ nicht in der Parameterform (8)
an, sondern mit $x$ als unabhängiger Variabeln
\[
J=\int f(x,y,y')\,dx,\quad y'=\frac{dy}{dx},
\]
wo dann in den Formeln (24) bis (27) die Funktionen $F_{x'},F_{y'}$ bezw. durch
\[
f(x,y,p)-pf_{y'}(x,y,p),\quad f_{y'}(x,y,p)
\]
zu ersetzen sind, wobei nunmehr $p$ die Gefällfunktion des Feldes bedeutet.} zuerst die partielle Differentialgleichung
(25) und schliesst daraus auf die Existenz einer Funktion
$V(x,y)$, für welche die Gleichungen (24) und die partielle Differentialgleichung
(26) bestehen. Diese Funktion $V$ schreibt er in der Form des Linienintegrals
(27); ihre geometrische Bedeutung als Feldintegral bleibt ihm jedoch für den
allgemeinen Fall verborgen, da ihm der Begriff der Transversale noch fehlt.
Diese allgemeinen Resultate wendet er dann auf den Fall der geodätischen
Linien an. Er erhält so die Gleichungen (29) und (30) und zeigt nun für
diesen speziellen Fall, dass hier die der Funktion $V$ entsprechende Funktion
die Bedeutung des senkrechten geodätischen Abstandes von einer orthogonalen
Trajektorie der betrachteten Schar von geodätischen Linien besitzt.

Diese Arbeit von \textsc{Beltrami} scheint von Seiten der Variationsrechnung
gänzlich unbeachtet geblieben zu sein. Erst mehr als dreissig Jahre später
hat \textsc{Hilbert}\footnote{D. \textsc{Hilbert}, \textit{Mathematische Probleme}, Göttinger Nachrichten 1900, S. 292 und Archiv der Mathematik
und Physik (3), Bd. 1 (1901), S. 231—236.} dieselben Resultate aufs neue entdeckt und darauf seinen
\clearpage\SourceLeaf{520}
\ArithmeticHead{63}{R}{BEZIEHUNGEN ZUM UNABHÄNGIGKEITSSATZ.}
\noindent bekannten Hinlänglichkeitsbeweis für ein Extremum des Integrals $J$ gegründet.
Aus der Darstellung von \textsc{Beltrami} geht nicht unzweideutig hervor, ob er den
Anstoss zu seiner Untersuchung von \textsc{Gauss} erhalten hat; doch darf man dies
wohl als wahrscheinlich annehmen, und in diesem Fall würde sich nicht nur
ein innerer logischer, sondern auch ein äusserer historischer Zusammenhang
zwischen den \textsc{Gauss}schen Untersuchungen und den Fundamentalsätzen der
modernen Variationsrechnung ergeben.
\par\noindent\hspace*{\parindent}c) Verwandte Untersuchungen von \textsc{Hamilton} und weitere Verallgemeinerungen.
\par\noindent\hspace*{\parindent}12. Mit den hier besprochenen Sätzen »verwandt« sind Untersuchungen
über \textls[100]{Strahlensysteme}, die \textsc{Hamilton}\footnote{W. R. \textsc{Hamilton}, loc. cit. 9), Art. 8, 9, 19, 20.} fast gleichzeitig mit den \textsc{Gauss}schen
\textit{Disquisitiones} veröffentlicht hat, und die sich auf das spezielle Variationsproblem
\[
\int\sqrt{x'^2+y'^2+z'^2}\,dt=\text{Minimum}
\]
beziehen. Eine gegenseitige Beeinflussung zwischen \textsc{Gauss} und \textsc{Hamilton} ist
dabei ausgeschlossen, da die Arbeit von \textsc{Hamilton} vom 3. Dezember 1824, die
\textsc{Gauss}sche vom 8. Oktober 1827 datiert ist und beide Arbeiten im Laufe des
Jahres 1828 erschienen sind.

\textsc{Hamilton} betrachtet ein Strahlensystem, von dem stillschweigend vorausgesetzt
wird, dass es ein Feld bildet, und bezeichnet mit $\alpha(x,y,z),\beta(x,y,z)$,
$\gamma(x,y,z)$ die Richtungskosinus des Feldes, d. h. die Richtungskosinus des durch
den Punkt $x,y,z$ gehenden Strahls. Er beweist dann den Fundamentalsatz:
»Soll das betrachtete Strahlensystem ein \textls[100]{Normalensystem} bilden, so muss
der Differentialausdruck $\alpha(x,y,z)\,dx+\beta(x,y,z)\,dy+\gamma(x,y,z)\,dz$ das vollständige
Differential einer Funktion $V(x,y,z)$ sein, die \textsc{Hamilton} die \textls[100]{charakteristische Funktion} des Normalensystems nennt, sodass also
\begin{equation*}\tag{31}
\alpha=\frac{\partial V}{\partial x},\quad\beta=\frac{\partial V}{\partial y},\quad\gamma=\frac{\partial V}{\partial z}.
\end{equation*}

Ist diese Bedingung erfüllt, so schneiden die Flächen $V=\text{const.}$ die
Strahlen des Systems senkrecht. Die Funktion $V$ ist nichts anderes als der
senkrechte Abstand $\rho$ des Punktes $x,y,z$ von der Fläche $V=0$ dieser Schar.
Endlich genügt die Funktion $V$ der partiellen Differentialgleichung
\clearpage\SourceLeaf{521}
\ArithmeticHead{64}{L}{BOLZA, GAUSS UND DIE VARIATIONSRECHNUNG.}
\begin{equation*}\tag{32}
\left(\frac{\partial V}{\partial x}\right)^2+\left(\frac{\partial V}{\partial y}\right)^2+\left(\frac{\partial V}{\partial z}\right)^2=1.
\end{equation*}
Weiter beweist \textsc{Hamilton} folgende partielle Differentialgleichungen für die
Funktionen $\alpha,\beta,\gamma$
\begin{equation*}\tag{33}
\left\{\begin{array}{@{}r@{}l@{}l@{}l@{}}
&\beta\left(\frac{\partial\alpha}{\partial y}-\frac{\partial\beta}{\partial x}\right)&+\gamma\left(\frac{\partial\alpha}{\partial z}-\frac{\partial\gamma}{\partial x}\right)&=0,\\[1.1ex]
\alpha\left(\frac{\partial\beta}{\partial x}-\frac{\partial\alpha}{\partial y}\right)&&+\gamma\left(\frac{\partial\beta}{\partial z}-\frac{\partial\gamma}{\partial y}\right)&=0,\\[1.1ex]
\alpha\left(\frac{\partial\gamma}{\partial x}-\frac{\partial\alpha}{\partial z}\right)&+\beta\left(\frac{\partial\gamma}{\partial y}-\frac{\partial\beta}{\partial z}\right)&&=0,
\end{array}\right.
\end{equation*}
aus denen sich ein zweiter Beweis des Fundamentalsatzes ergibt.

In den späteren Teilen seiner grossen Arbeit über Strahlensysteme hat
\textsc{Hamilton}\footnote{loc. cit. 9), First supplement, Bd. 16, Part I (1830); Second supplement, Bd. 16, Part II (1831,
datiert 1830); Third supplement, Bd. 17 (1837, datiert 1832).} seine Untersuchungen auf den weit allgemeineren Fall des Extremums
des Integrals
\[
\int n\left(x,y,z,\frac{dx}{ds},\frac{dy}{ds},\frac{dz}{ds}\right)\,ds
\]
ausgedehnt\footnote{Hierauf macht G. \textsc{Prange} aufmerksam in seiner Habilitationsrede, W. R. \textsc{Hamilton}s \textit{Bedeutung
für die geometrische Optik}, Jahresbericht der Deutschen Mathematikervereinigung, Bd. 30 (1921), S. 72.
Vgl. dazu auch die Arbeit von G. \textsc{Prange}, \textit{W. R. Hamiltons Arbeiten zur Optik und Mechanik. Ihre
Bedeutung und ihr Einfluss auf die Entwicklung dieser Wissenschaften}, die demnächst in den Acta der
Leopoldinisch-Carolinischen Akademie der Naturforscher erscheinen wird.}, d. h. des Integrals (19) für $n=3$.
\par\clearpage\endgroup
